\section{A concrete quadratic character from the eleventh multiplication formula}
\label{sec:odd-anderson-characters}

This section constructs the explicit continuous quadratic character
$\Theta_{17,11}$, which vanishes on the arithmetic subgroup.  Its
nonvanishing on the full
profinite Grothendieck--Teichm\"uller group is not established.  A separate
formal example shows why beta factorization, the unitary relation, and
the dyadic multiplication formulas do not by themselves eliminate these
characters.  The example is not asserted to lift to a pentagon element.

Let $G_0$ be the profinite gentle group and let
$G=\widehat{\mathrm{GT}}\subset G_0$.  Write
$F=\widehat F(x,y)$ and $z=(xy)^{-1}$.  We use the convention
\[
 \phi_g(x)=x^{\lambda_g},\qquad
 \phi_g(y)=f_g^{-1}y^{\lambda_g}f_g,
 \qquad f_g\in F'.
\]
Let $B_g(X,Y)$ be the normalized Alexander coefficient, so that the
semilinear coefficient on $F'/F''$ is
\[
 A_g(X,Y)=S_{\lambda_g}(X)S_{\lambda_g}(Y)B_g(X,Y),
 \qquad S_a(T)=\frac{T^a-1}{T-1}.
\]
Composition gives
\begin{equation}
 B_{gh}(X,Y)=B_g(X,Y)
                 B_h(X^{\lambda_g},Y^{\lambda_g}).
 \label{eq:odd-beta-cocycle}
\end{equation}
These are the conventions of \cite[\S2.3, p.~101]{Furusho2003}.

\subsection{The boundary characters in the unitary relation}

\begin{lemma}\label{lem:all-finite-unitarity}
Suppose that $\lambda_g=1$.  Let $p$ be an odd prime and let
$a,b\in\overline{\mathbb F}_p^\times$ have finite orders prime to $p$.
Then
\[
 B_g(a,b)B_g(a^{-1},b^{-1})=1.
\]
In fact $B_g(a,b)=1$ if $a=1$, $b=1$, or $ab=1$.
\end{lemma}

\begin{proof}
Choose $m$ prime to $p$ divisible by the orders of $a$ and $b$, and
put $U_m=\ker(F\to(\mathbb Z/m)^2)$.  Since $f_g\in F'$, the
automorphism $\phi_g$ induces the identity on $F/\gamma_3F$.  Choose
$c\in F$ such that $\phi_g(z)=c^{-1}zc$, and write
$[c]=(u,v)\in F^{\mathrm{ab}}=\widehat{\mathbb Z}^2$.  In the
class-two quotient the commutator subgroup is freely generated over
$\widehat{\mathbb Z}$ by $[x,y]$.  The equality $c^{-1}zc=z$ in this
quotient implies $u=v$.

The cover with fundamental group $U_m$ has deck group $(\mathbb Z/m)^2$.
Its boundary subrepresentation in $H_1(U_m,\mathbb F_p)$ is the sum
of the three permutation representations induced from the cyclic
subgroups generated by the images of $x$, $y$, and $z$, with one
trivial relation among the boundary classes.  On the boundaries above
$x$ the conjugator is $1$; above $y$ it is $f_g$, whose deck image is
$1$; above $z$ it is $c$, whose deck image is diagonal.  Thus $\phi_g$
fixes the boundary representation pointwise.  More explicitly, a
character $(a,b)$ occurring above $z$ satisfies $ab=1$, and its
conjugator scalar is $a^u b^v=(ab)^u=1$.  The corresponding assertions
above $x$ and $y$ are immediate.  Every nontrivial deck character
occurs in $H_1(U_m,\overline{\mathbb F}_p)$ with multiplicity one.
The coefficient $B_g(a,b)$ is its scalar action.  This proves the
last assertion for all nontrivial boundary characters.  At the trivial
character $B_g(1,1)=1$, since the class-two commutator is fixed.

For completeness, the complement of the boundary subrepresentation
carries the usual nondegenerate alternating form.  One can see its
preservation without a topological realization of $\phi_g$.  A surface
presentation of $U_m$ has the form
\[
 \prod_{j=1}^{(m-1)(m-2)/2}[a_j,b_j]
       \prod_{i=1}^{3m}c_i=1.
\]
In the class-two exponent-$p$ quotient, project the degree-two
commutators to the exterior square of homology modulo the boundary
classes.  The relation yields the bivector
$\sum_j[a_j]\wedge[b_j]$.  Conjugating a boundary word changes its
degree-two coordinate by a wedge containing its boundary class, so
this change vanishes after projection.  Positive boundary preservation
therefore fixes the bivector and the associated alternating form.

The deck action preserves the form, so the character spaces at
$(a,b)$ and $(a^{-1},b^{-1})$ pair perfectly whenever
$a,b,ab\ne1$.  Preservation of the form makes the product of their
scalar actions equal to one.  The boundary calculation proves the
same identity in the remaining cases.
\end{proof}

\subsection{An explicit family of arithmetic-trivial characters}

\begin{theorem}\label{thm:odd-anderson-character}
Let $p\equiv5\pmod{12}$ be prime, and let $n\geq1$ satisfy
$\gcd(n,3p)=1$.  Choose a primitive cube root
$\omega\in\mathbb F_{p^2}$.  The expression
\begin{equation}
 E_{p,n}(g)=
 \frac{\displaystyle\prod_{k=0}^{n-1}B_g(\omega,\omega^k)}
      {\displaystyle\prod_{\xi^n=1}B_g(\xi,\omega)}
 \label{eq:odd-anderson-defect}
\end{equation}
belongs to the norm-one subgroup of $\mathbb F_{p^2}^\times$.
Here all evaluations are reduced modulo $p$, and the denominator
ranges over all $n$th roots in $\overline{\mathbb F}_p$.
The rule
\begin{equation}
 (-1)^{\Theta_{p,n}(g)}=E_{p,n}(g)^{(p+1)/2}
 \label{eq:odd-anderson-sign}
\end{equation}
defines a continuous homomorphism $G_0\to\mathbb F_2$, independent
of the choice between $\omega$ and $\omega^2$.  It vanishes on
$G_{\mathbb Q}$.

In particular its restriction to $G$ is an explicit candidate for a
nonzero element of
$\operatorname{Hom}_{\mathrm{cts}}(G/\langle\!\langle
G_{\mathbb Q}\rangle\!\rangle,\mathbb F_2)$.
The theorem does not assert that this restriction is nonzero.
\end{theorem}

\begin{proof}
The evaluations are nonzero because they are coefficients of
invertible semilinear actions on the appropriate character spaces.
The $p^2$-power Frobenius fixes $\omega$ and permutes the $n$th
roots.  Hence the expression lies in $\mathbb F_{p^2}^\times$.
Equation~\eqref{eq:odd-beta-cocycle} and reindexing the denominator
give
\begin{equation}
 E_{p,n}(gh)=E_{p,n}(g)\,
             \sigma_{\lambda_g}\bigl(E_{p,n}(h)\bigr),
 \label{eq:odd-anderson-crossed}
\end{equation}
where $\sigma_\lambda$ is the identity for $\lambda\equiv1\pmod3$
and the $p$-power Frobenius for $\lambda\equiv-1\pmod3$.
Indeed $\lambda$ permutes all $n$th roots, and its only remaining
effect in the displayed expression is its action on $\omega$.

If $\lambda_g=1$, Lemma~\ref{lem:all-finite-unitarity} gives
$E_{p,n}(g)^p=E_{p,n}(g)^{-1}$.  In the denominator one reindexes
the $p$th powers of the $n$th roots as their inverses; in the numerator
$\omega^p=\omega^{-1}$.  For arbitrary $g$, choose
$h\in G_{\mathbb Q}$ with $\lambda_h=\lambda_g$ and write $g=uh$,
where $\lambda_u=1$.  The arithmetic Anderson multiplication formula
\cite[Corollary~8.6.3]{Anderson1989}, in the beta form stated explicitly
in \cite[Remark~2.4, p.~101]{Furusho2003}, gives
\begin{equation}
 E_{p,n}(h)=\omega^{n\Psi_n(h)}\in\mu_3.
 \label{eq:odd-anderson-arithmetic}
\end{equation}
Its exact sign convention is immaterial here; in the convention above
$\Psi_n=-\rho_n$.  Since $p\equiv-1\pmod3$, the group $\mu_3$ also
has norm one.  Equation~\eqref{eq:odd-anderson-crossed} now proves
that $E_{p,n}(g)$ has norm one.

The norm-one subgroup is cyclic of order $p+1$, so
\eqref{eq:odd-anderson-sign} is a sign.  Frobenius acts on that
subgroup by inversion and therefore fixes its quadratic sign.
Consequently \eqref{eq:odd-anderson-crossed} makes $\Theta_{p,n}$
an ordinary homomorphism.  It is continuous because it uses finitely
many finite character evaluations.  Since $3$ divides $(p+1)/2$,
\eqref{eq:odd-anderson-arithmetic} makes this homomorphism zero on
$G_{\mathbb Q}$.  Replacing $\omega$ by $\omega^2$ applies Frobenius
and leaves the sign unchanged.
\end{proof}

For a pentagon element, Ihara's beta factorization gives
$B_g(s,t)=\Gamma_g(s)\Gamma_g(t)/\Gamma_g(s+t)$ in additive
notation.  Consequently
\begin{equation}
 E_{p,n}(g)=
 \frac{\displaystyle\prod_{c\in\frac1n\mathbb Z/\mathbb Z}
                    \Gamma_g(\tfrac13+c)}
      {\displaystyle\Gamma_g(\tfrac n3)
          \prod_{c\in\frac1n\mathbb Z/\mathbb Z}\Gamma_g(c)}.
 \label{eq:odd-anderson-gamma}
\end{equation}
This expression is independent of the character ambiguity in
$\Gamma_g$.  If the Anderson relation $(A_n)$ holds for $g$, then
$\Theta_{p,n}(g)=0$.  Thus the known dyadic multiplication formulas
eliminate the cases $n=2^a$.  Enriquez's theorem that
$\mathrm{GTK}=\widehat{\mathrm{GT}}$ also eliminates every proposed
character obtained from a multiplicative defect of the integral
Kummer coordinates: his Remark~8.6 gives
$\rho_{ab}=\rho_a+\rho_b$ before any reduction or division.
The exact reference is \cite[Theorem~8.1 and Remark~8.6]{Enriquez2007}.
It does not assert the equality
$\mathrm{GTA}=\widehat{\mathrm{GT}}$; the relation $(K_n)$ is recorded
as implying the logarithmic derivative of $(A_n)$ in
\cite[p.~102]{Furusho2003}.

\begin{corollary}\label{cor:concrete-radical-anderson-character}
Let $H=G_{\mathbb Q}$ and let $N_{\mathrm{rad}}$ be the simultaneous
kernel in $G$ of the multiplier and the rational-prime Kummer cocycles.
Put
\[
 L_{\mathrm{rad}}=\mathbb Q(\mu_\infty,q^{1/m}:
                          q\text{ prime},\ m\geq1),
 \qquad K_{\mathrm{rad}}=G_{L_{\mathrm{rad}}}.
\]
Assume the simultaneous normalization
$G=N_{\mathrm{rad}}H$, with
$N_{\mathrm{rad}}\cap H=K_{\mathrm{rad}}$, and identify the relative
abelian quotient as
\[
 B=N_{\mathrm{rad}}/
       \overline{[G,N_{\mathrm{rad}}]K_{\mathrm{rad}}}.
\]
For a primitive cube root $\omega\in\mathbb F_{289}$ the formula
\begin{equation}
 (-1)^{\psi_{17,11}(u)}=
    \left(\prod_{\substack{\xi^{11}=1\\\xi\ne1}}
                        B_u(\xi,\omega)\right)^9,
 \qquad u\in N_{\mathrm{rad}},
 \label{eq:concrete-radical-anderson-sign}
\end{equation}
defines a continuous $G$-invariant quadratic character vanishing on
$K_{\mathrm{rad}}$.  It therefore gives a precisely specified element
of $\operatorname{Hom}_{\mathrm{cts}}(B,\mathbb F_2)$.
Its unique extension to $G$ vanishing on $H$ is $\Theta_{17,11}$.
In particular $\psi_{17,11}$ is nonzero if and only if
$\Theta_{17,11}|_G$ is nonzero; neither nonvanishing assertion is
proved here.
\end{corollary}

\begin{proof}
Take $p=17$, $n=11$ in Theorem~\ref{thm:odd-anderson-character}.
For $u\in N_{\mathrm{rad}}$ the multiplier is one.
In the numerator of \eqref{eq:odd-anderson-defect}, four values of
$k$ are $1$ modulo $3$, four are $0$ modulo $3$, and three are $2$
modulo $3$.  Lemma~\ref{lem:all-finite-unitarity} gives
\[
 \prod_{k=0}^{10}B_u(\omega,\omega^k)=B_u(\omega,\omega)^4,
 \qquad B_u(\omega,\omega)^{18}=1.
\]
Its ninth-power contribution is therefore one.  The denominator's
$\xi=1$ term is also one.  Since inversion fixes a sign, the
restriction of \eqref{eq:odd-anderson-sign} is exactly
\eqref{eq:concrete-radical-anderson-sign}.

Being a restriction of an ordinary character of $G$, this map is
$G$-invariant; being zero on $H$, it is zero on $K_{\mathrm{rad}}$.
In fact the arithmetic formula gives $E_{17,11}(u)=1$ on
$K_{\mathrm{rad}}$, since $\Psi_{11}(u)=0$ there.
It consequently factors through the displayed quotient $B$.
The factorization $G=N_{\mathrm{rad}}H$ proves both uniqueness
of the extension and the equivalence of the two nonvanishing claims.
\end{proof}

There is also a finite-dimensional interpretation of the product in
\eqref{eq:concrete-radical-anderson-sign}.  Take the degree-$33$ cover
with deck group $C_{11}\times C_3$, given algebraically by
$u^{11}+v^3=1$ and the map $t=u^{11}$.  There are three boundary
components above $t=0$, eleven above $t=1$, and one above infinity.
Euler characteristic therefore gives genus ten for its compactification.
Its homology has the ten character lines $(\xi,\omega)$ with
$\xi$ a nontrivial $11$th root, together with their inverse lines.
The individual evaluations may lie in $\mathbb F_{17^{10}}$;
$\mu_{11}$ is not contained in $\mathbb F_{289}$.
The sum descends to a ten-dimensional space over $\mathbb F_{289}$,
since the $17^2$-power Frobenius permutes these ten characters.  It is
the whole $\omega$-eigenspace of the cubic deck action on the compact
homology.  The product is the determinant of an element of
$N_{\mathrm{rad}}$ on this space, and
\eqref{eq:concrete-radical-anderson-sign} is its quadratic norm-one
sign.  Thus the proposed test involves one finite homology
representation of a genus-ten curve.

\subsection{A formal sign ambiguity which survives duplication}

\begin{proposition}\label{prop:formal-odd-sign}
There exists a Frobenius-equivariant sign-valued formal gamma function
at the prime $17$ whose beta coboundary satisfies beta associativity,
symmetry, the unitary relation, and the duplication formula, but for
which the expression in \eqref{eq:odd-anderson-gamma} with $n=11$
equals $-1$.  It may be taken constant in the pro-$17$ directions,
so all its logarithmic derivatives vanish.

This is a non-implication statement for the listed formal identities.
It does not provide a solution of the full nonabelian relations
$(K_n)$ or of the pentagon equation.
\end{proposition}

\begin{proof}
Let $T=(\mathbb Q/\mathbb Z)_{17'}$ be the subgroup of elements of
order prime to $17$.  In $(\mathbb Z/33)^\times$ put
\[
 S=\langle2\rangle
   =\{1,2,4,8,16,32,31,29,25,17\}.
\]
Define $\gamma:T\to\mathbb F_2$ by
\[
 \gamma(s)=
 \begin{cases}
  1,&s=a/33\text{ for }a\in S,\\
  0,&\text{otherwise},
 \end{cases}
 \qquad \Gamma(s)=(-1)^{\gamma(s)}.
\]
For an arbitrary rational torsion argument $s$, the definition means
$\gamma(s)=\gamma(s_{17'})$, where $s_{17'}$ is its prime-to-$17$
component.  Thus the value is constant in the $17$-primary variable;
it is not declared to be $1$ merely because the denominator is
divisible by $17$.  In particular $\gamma$ is zero whenever the
prime-to-$17$ component has even order.
The set $S$ is stable under $2$, $-1=2^5$, and $17=2^9$.
Thus $\Gamma(-s)=\Gamma(s)$ and $\Gamma(17s)=\Gamma(s)$.

We check the duplication relation
\begin{equation}
 \gamma(s)+\gamma(s+\tfrac12)=\gamma(2s)+\gamma(\tfrac12).
 \label{eq:formal-sign-duplication}
\end{equation}
The last term is zero.  If $s$ has odd order, then $s+1/2$ has
even order and the assertion is the invariance under doubling on
odd torsion.  If the $2$-primary component of $s$ has order exactly
two, put $v=s+1/2$.  Then $v$ has odd order, $2s=2v$, and the same
invariance proves the identity.  If the $2$-primary component has
order at least four, all three possibly nonzero arguments on the
left and right have even order, so all terms vanish.  This proves
\eqref{eq:formal-sign-duplication}, or equivalently
\[
 \Gamma(s)\Gamma(s+\tfrac12)=\Gamma(2s)\Gamma(\tfrac12).
\]
Iteration gives all dyadic multiplication identities with zero Kummer
correction.

Put $B(s,t)=\Gamma(s)\Gamma(t)/\Gamma(s+t)$.  It is symmetric,
satisfies beta associativity, is one on the axes and on $s+t=0$,
and satisfies $B(s,t)B(-s,-t)=1$.  These assertions follow directly
from the coboundary definition and the fact that all values are signs.
For the prime-to-$17$ spectral components, the Frobenius invariance
allows the assignment to descend to $\mathbb Z_{17}$ coefficients.
Make it constant on each pro-$17$ formal neighborhood.  Equivalently,
use the product decomposition of
$\mathbb Z_{17}[[\widehat{\mathbb Z}]]$ into the unramified formal
power-series rings indexed by Frobenius orbits of prime-to-$17$
characters.  The specified constants define a unit in this completed
group algebra.  The beta coboundary is a corresponding unit in the
two-variable algebra, and its logarithmic derivatives are zero.

Finally $\Gamma(11/3)=\Gamma(2/3)=1$, and $\Gamma(c)=1$ for
every $c\in\frac1{11}\mathbb Z/\mathbb Z$.  The numerators of the
points $1/3+k/11$, expressed with denominator $33$, are $11+3k$.
Their intersection with $S$ is
\[
 \{2,8,17,29,32\},
\]
which has five elements.  Therefore the numerator of
\eqref{eq:odd-anderson-gamma} is $(-1)^5=-1$ and its denominator
is $1$.  At $p=17$ the quadratic exponent is $(17+1)/2=9$, so
the associated formal quadratic sign is also $-1$.
\end{proof}

The remaining question is a full-group realization or elimination
problem.  Neither the finite sign computation nor the fact that
logarithmic derivatives forget this sign permits one to infer an
element of $\widehat{\mathrm{GT}}$ with nonzero
$\Theta_{17,11}$.  Conversely, eliminating these candidates requires
an additional consequence of the full relations, rather than only
the four formal identities in Proposition~\ref{prop:formal-odd-sign}.

\subsection{Two limitations of a direct eleventh-Kummer comparison}

Put $D=\ker\lambda\cap\bigcap_q\ker\rho_q^I$ in the full
profinite Grothendieck--Teichm\"uller group.  By \cite[Theorem~8.1 and Remark~8.6]{Enriquez2007}, the exact
Kummer relation, with zero Kummer coordinate on $D$, is
\begin{equation}\label{eq:eleventh-local-system-K}
d_n(f_g)=f_g\qquad(g\in D).
\end{equation}
Here $H_n\leq\widehat F(x,y)$ is the kernel of $x\mapsto1,y\mapsto0$
in $\mathbb Z/n$, freely generated by
\[
 X=x^n,\qquad Y_j=x^{-j}yx^j\quad(0\leq j<n),
\]
and $d_n:H_n\to\widehat F(x,y)$ sends $X\mapsto x$, $Y_0\mapsto y$
and $Y_j\mapsto1$ for $j\ne0$.  We apply the full nonabelian relation \eqref{eq:eleventh-local-system-K}.

\begin{lemma}\label{lem:kn-different-local-systems}
For $n=11$ and a primitive cube root $\omega$ in
$\mathbb F_{17^2}$, consider characters over $\overline{\mathbb F}_{17}$.
The character of $H_{11}$ obtained by pulling
back $x\mapsto1,y\mapsto\omega$ along $d_{11}$ is different from
the restriction of every character
$x\mapsto\xi,y\mapsto\omega$ with $\xi^{11}=1$.
Nor is either kind a deck translate of the other.
\end{lemma}
\begin{proof}
The pullback along $d_{11}$ has values
\[
 X\mapsto1,\qquad Y_0\mapsto\omega,\qquad
 Y_j\mapsto1\ (j\ne0).
\]
Each character of the second kind has values
\[
 X\mapsto\xi^{11}=1,\qquad Y_j\mapsto\omega
 \quad\hbox{for every }j.
\]
The deck generator cyclically permutes the conjugacy classes of
the $Y_j$.  It therefore permutes the eleven characters having a
single nontrivial $Y_j$-value and fixes the character whose eleven
$Y_j$-values are all $\omega$.  These two orbits are disjoint.
\end{proof}

Consequently the direct operation ``apply \eqref{eq:eleventh-local-system-K}, then pass to the
pulled-back rank-one local system'' does not identify the product
\begin{equation}\label{eq:eleventh-local-system-P}
P_g=\prod_{\substack{\xi^{11}=1\\\xi\ne1}}B_g(\xi,\omega)
\end{equation}
with the determinant on that pulled-back local system.  This
observation is a limitation of this proposed argument, not a claim
that the full relation \eqref{eq:eleventh-local-system-K} cannot constrain \eqref{eq:eleventh-local-system-P} by another route.

Here is a separate, explicit limitation of a pairing-only argument.
Let $C$ be the smooth compact curve with affine equation
$u^{11}+v^3=1$, and let
$A=C_{11}\times C_3$ be its deck group over the $t$-line,
$t=u^{11}$.  We use the geometric finite homology description:
over an algebraic closure of $\mathbb F_{17}$, its $20$-dimensional
homology has exactly the character lines
\begin{equation}\label{eq:eleventh-local-system-H}
(\xi^a,\omega^b),\qquad
 a\in(\mathbb Z/11)^\times,\quad b\in\{1,2\}.
\end{equation}
Every line is one-dimensional.  The alternating intersection form
pairs a line only with its inverse-character line.
For completeness, the character description follows from the
two-loop cellular complex of the three-punctured sphere.  On a
nontrivial character the boundary from the two one-cells to the
one zero-cell has rank one, so the open cover has a single
homology line.  Compactification kills precisely the peripheral
lines whose character is trivial on one of the three inertia
groups.  The inertia values here are $\xi^a$, $\omega^b$ and
$(\xi^a\omega^b)^{-1}$.  Since eleven and three are coprime,
all three are nontrivial exactly for the pairs in \eqref{eq:eleventh-local-system-H}.
The trivial character contributes no compact homology.

\begin{proposition}\label{prop:kn-linear-sign-model}
There is an $\mathbb F_{17}$-linear automorphism of the space \eqref{eq:eleventh-local-system-H}
that commutes with the whole deck group $A$, preserves its
alternating intersection form, and has determinant $-1$ on the
ten-dimensional $\omega$-eigenspace over $\mathbb F_{17^2}$.
Its quadratic norm-one determinant sign is therefore $-1$.
This is a linear model, not a GT element or a solution of the
pentagon or of \eqref{eq:eleventh-local-system-K}.
\end{proposition}
\begin{proof}
The square of Frobenius acts on the nonzero exponents modulo
eleven by multiplication by $17^2\equiv3$.  Its two orbits are
\[
 A_0=\{1,3,9,5,4\},\qquad
 A_1=\{2,6,7,10,8\}=-A_0.
\]
Ordinary Frobenius acts on the pairs $(a,b)$ in \eqref{eq:eleventh-local-system-H} as
\[
 (a,b)\longmapsto(6a,-b).
\]
It thus preserves the set
\[
 \mathcal O=(A_0\times\{1\})\cup(A_1\times\{2\}).
\]
Indeed $6A_0=A_1$ and $6A_1=A_0$.
The set $\mathcal O$ is also stable under inversion
$(a,b)\mapsto(-a,-b)$, as is its complement.

Define the automorphism to be $-1$ on the ten character lines in
$\mathcal O$ and $+1$ on the ten remaining lines.  Frobenius
stability makes it descend to $\mathbb F_{17}$, and its scalar
action on each character line makes it commute with $A$.
The intersection form pairs only mutually inverse characters:
if $v,w$ have characters $\chi,\psi$, invariance gives
$\langle v,w\rangle=\chi(a)\psi(a)\langle v,w\rangle$ for every
deck element $a$, which forces $\psi=\chi^{-1}$ whenever the
pairing is nonzero.  Our automorphism has equal signs on these
paired lines, so it preserves that form.

There are exactly five negative lines in the $\omega$-eigenspace,
namely those indexed by $A_0$.  Its determinant there is
$(-1)^5=-1$.  This value belongs to the norm-one group of
$\mathbb F_{17^2}^\times$, and raising it to
$(17+1)/2=9$ leaves $-1$, as asserted.
\end{proof}

The exact remaining question is whether the full nonabelian
relations, including \eqref{eq:eleventh-local-system-K} and the pentagon, exclude this sign in
the GT action.  Neither Lemma~\ref{lem:kn-different-local-systems}
nor Proposition~\ref{prop:kn-linear-sign-model} answers it.
In particular, no assertion about nonvanishing of
$\Theta_{17,11}|_{\widehat{\mathrm{GT}}}$ follows from these two statements.

